% Kompiuterijos katedros šablonas
% Template of Department of Computer Science II
% Versija 1.0 2015 m. kovas [ March, 2015]

\documentclass[a4paper,12pt,fleqn]{article}
\usepackage[unicode,colorlinks=false]{hyperref}


\usepackage[utf8x]{inputenc}
%

\usepackage[L7x]{fontenc}
\usepackage{times}
\usepackage{ucs}

 %package to switch the language
\usepackage{etoolbox}

  %set up of the page margins
\usepackage[top=2cm, bottom=2cm, left=3cm, right=1.5cm]{geometry}

 %1.1 line spacing
\linespread{1.1}


  %page numbering at the right side
\usepackage{fancyhdr}
\pagestyle{fancyplain}
\fancyhf{}
\renewcommand{\headrulewidth}{0pt} 
\fancyhfoffset[RO]{0cm}

  %to number at the bottom (exchange lines to number at the top)
\rfoot{\thepage}
  %\rhead{\thepage} %

% \usepackage[usenames,dvipsnames]{pstricks}
\urlstyle{same}
\hypersetup{
%  citecolor=Blue,
%  linkcolor=Blue,
%  urlcolor=Blue
pdfborder={0 0 0 }
}

 %for includegraphics
\usepackage{graphicx}



\usepackage[toc,page]{appendix}


\usepackage{caption}

 %for source codes
\usepackage{listings}
\lstset{commentstyle=\color{red},xleftmargin=10pt, framexleftmargin=6pt, numbersep=1mm, frame=single, numbers=left,numberstyle=\footnotesize,extendedchars=\true, inputencoding=utf8x,basicstyle=\footnotesize,extendedchars=true,
 keywordstyle=\color{black}\bfseries, breaklines=true, breakautoindent=true,framesep=8pt,linewidth=0.95\textwidth
}

 %for algorithms
\usepackage{algorithm}
\usepackage{algorithmic}
 %instead of the above two packages we can use algorithms2e
 %\usepackage[boxed,linesnumbered,vlined,slide]{algorithm2e}

 %special symbols
\usepackage{amsfonts}
\usepackage{amssymb}
\usepackage{amsmath}

 %for theorem like environments
\usepackage{amsthm}

 \usepackage{datetime}
 \renewcommand{\dateseparator}{--}


% SI system units
\usepackage{siunitx}
\sisetup{detect-all}
% Problem with fonts \SI{x.xx}{\micro\metre}, solved with updmap-sys --enable Map=utm.map
\renewcommand{\sfdefault}{uhv}
\renewcommand{\rmdefault}{utm}
\renewcommand{\ttdefault}{ucr}

% List management (itemize, etc.)
\usepackage{enumitem}

\newcommand*{\urlw}[1]{\href{#1}%
            {\nolinkurl{#1}}}

\numberwithin{equation}{section}


\newtoggle{inLithuanian}
 %If the report is in Lithuanian, it is set to true; otherwise, change to false
\settoggle{inLithuanian}{false}

%create file preface.tex for the preface text
%if preface is needed set to true
\newtoggle{needPreface}
\settoggle{needPreface}{false}

\newtoggle{signaturesOnTitlePage}
\settoggle{signaturesOnTitlePage}{true}


\theoremstyle{definition}
\newtheorem{definition}{\keyWordDefinition}
\newtheorem{example}{\keyWordExample}
\def\QED{\unskip\nobreak\hfill\kern5pt$\Box$}
\def\code#1{\texttt{#1}}

\iftoggle{inLithuanian}{
%\usepackage[L7x]{fontenc}
\usepackage[english,lithuanian]{babel}

\newcommand{\todayiso}{\the\year \dateseparator \twodigit\month \dateseparator \twodigit\day}


\renewcommand{\today}{\number\year\space m. \space \ifcase\month\or
  sausio\or vasario\or kovo\or balandžio\or gegužės\or birželio\or
  liepos\or rugpjūčio\or rugsėjo\or spalio\or lapkričio\or
  gruodžio\fi
  \space\number\day\space d.}


 \usepackage{tocloft}
 \renewcommand\cftsecaftersnum{.} 
 \renewcommand\cftsubsecaftersnum{.} 
 \renewcommand\cftsubsubsecaftersnum{.}

 \usepackage{VUMIFKK}

 \DeclareCaptionLabelFormat{captionlt}{#2 #1}
   %smth is not fine with algorithms 
 \DeclareCaptionLabelFormat{captionltalg}{#2 #1 algoritmas}

 \usepackage{indentfirst}
 \renewcommand{\appendixtocname}{Priedai}
 \renewcommand{\appendixpagename}{Priedai}
 \renewcommand{\contentsname}{Turinys} 

 \renewcommand{\lstlistingname}{išeities kodas}
 \renewcommand{\figurename}{pav}
 \renewcommand{\tablename}{lentelė}


 \captionsetup*[lstlisting]{   
 labelsep=period,labelformat=captionlt
 }
 \captionsetup*[figure]{   
% labelsep=period,
 labelsep=space, %babel redefines pav to pav.
 labelformat=captionlt
 }
 \captionsetup*[table]{   
  labelsep=period,
  labelformat=captionlt
 }
 \renewcommand{\algorithmicrequire}{\textbf{Įvestis:}}
 \renewcommand{\algorithmicensure}{\textbf{Išvestis:}}

 \captionsetup*[algorithm]{   
 labelsep=period,labelformat=captionltalg
 }

\renewcommand{\thmhead}[3]{#2 #1#3}

}
{
%\usepackage[OT1,T1]{fontenc}
%\usepackage[L7x]{fontenc}



\usepackage[english]{babel}
\newcommand{\todayiso}{\twodigit\month \dateseparator \twodigit\day \dateseparator \the\year}
 \captionsetup*[algorithm]{   
 labelsep=period
 }
\captionsetup*[lstlisting]{   
 labelsep=period
 }
 \captionsetup*[figure]{   
 labelsep=period
 }
 \captionsetup*[table]{   
 labelsep=period
 }


}

%some kywords
 \def\keywordAbstract{\iftoggle{inLithuanian}{Santrauka}{Abstract}}
 \def\keywordAbstractOther{\iftoggle{inLithuanian}{Summary}{Santrauka}}
 \def\keyWordIntroduction{\iftoggle{inLithuanian}{Įvadas}{Introduction}}
 \def\keyWordConclusions{\iftoggle{inLithuanian}{Išvados ir rekomendacijos}{Conclusions and Recommendations}}

 \def\keyWordPreface{\iftoggle{inLithuanian}{Pratarmė}{Preface}}
 \def\keyWordAppendice{\iftoggle{inLithuanian}{Priedas}{Appendix}}
 \def\keyWordSignature{\iftoggle{inLithuanian}{parašas}{signature}}
 \def\keyWordDefinition{\iftoggle{inLithuanian}{apibrėžimas}{Definition}}
 \def\keyWordExample{\iftoggle{inLithuanian}{pavyzdys}{Example}}

\newcommand{\bothabstracts}[3]{
\setcounter{secnumdepth}{0}
\newpage
\hspace{2cm}
{\centering{\section{\keywordAbstract}}}

#1
\newpage
\hspace{2cm}
{\centering \section{\keywordAbstractOther}}

\begin{center}{\textbf{#2} }\end{center}

 #3
\setcounter{secnumdepth}{3}
}

 %non-numbered sections: #1 param: for labeling sec:#1, #2 -section title
\newcommand{\sectionWithoutNumber}[2]{\newpage
%\hspace{2cm}
\section*{#1}
\label{sec:#2}
\addcontentsline{toc}{section}{\nameref{sec:#2}}%{#3}
 }



\newcommand{\referenceSources}[1]{
\newpage
\cleardoublepage
\phantomsection
\iftoggle{inLithuanian}{
 \renewcommand{\refname}{Literatūros šaltiniai}

 \addcontentsline{toc}{section}{Literatūros šaltiniai}
 \markboth{\refname}{Literatūros šaltiniai}
 }
{

\addcontentsline{toc}{section}{References}
\markboth{References}{References}
}

\bibliographystyle{plain}
\bibliography{#1}
}



 \newcommand\authorsignature[1]{
\begin{flushright}
 \begin{minipage}[b]{0.45\textwidth}
  \centering
  \rule{\textwidth}{0.5pt}\\
   #1
  \end{minipage}
\end{flushright}
 }




 \newcommand\authorsignatures[5]{%
   \vspace{1cm}
   \authorsignature{#1}
   \ifstrequal{#2}{}{}{\vspace{0.3cm}
     \authorsignature{#2}
     \ifstrequal{#3}{}{}{\vspace{0.3cm}
      \authorsignature{#3}
      \ifstrequal{#4}{}{}{\vspace{0.3cm}
        \authorsignature{#4}
        \ifstrequal{#5}{}{}{\vspace{0.3cm}
         \authorsignature{#5}       
        }
      }
    }
} 
}

\newcommand{\authortitle}{
\iftoggle{signaturesOnTitlePage}{
\tiny{\keyWordSignature}
}{}
}

\newcommand{\depttitlepage}[8]
{
\thispagestyle{empty}
\begin{center}

\includegraphics[width=2cm]{jb_VU_zenklas}

%\vspace{-1cm}

\iftoggle{inLithuanian}
{ 
  VILNIAUS UNIVERSITETAS\\
  MATEMATIKOS IR INFORMATIKOS FAKULTETAS\\
  INFORMATIKOS INSTITUTAS\\
  KOMPIUTERINIO IR DUOMENŲ MODELIAVIMO KATEDRA
}
{
  VILNIUS UNIVERSITY \\
  FACULTY OF MATHEMATICS AND INFORMATICS \\
  INSTITUTE OF COMPUTER SCIENCE\\
  DEPARTMENT OF COMPUTATIONAL AND DATA MODELING
}

\vspace{5cm}

#1\\
\vspace{0.5cm}
\textbf{\Large #2}
\end{center}

\vspace{5cm}


\hspace{0.5\textwidth}
\begin{minipage}{0.4\textwidth}
 \begin{flushleft} 
\iftoggle{inLithuanian}
{
 \ifstrequal{#3}{}{}{Atliko:\\[5pt]}
}
{
\ifstrequal{#3}{}{}{Done by:\\[5pt]}
}

%\noindent
\begin{tabular}{@{}lr}%\setlength\tabcolsep{0pt}
\ifstrequal{#3}{}{}{#3&\hspace{2cm}\authortitle\\[5pt]}
\ifstrequal{#4}{}{}{#4&\authortitle\\[5pt]}
\ifstrequal{#5}{}{}{#5&\authortitle\\[5pt]}
\ifstrequal{#6}{}{}{#6&\authortitle\\[5pt]}
\ifstrequal{#7}{}{}{#7&\authortitle\\}
\end{tabular}

\end{flushleft}

\end{minipage}

\vspace{0.5cm}
\hspace{0.5\textwidth}
\begin{minipage}{0.4\textwidth}
 \begin{flushleft} 

\ifstrequal{#8}{}{}
{

\iftoggle{inLithuanian}
{
Vadovas:
}
{
Supervisor:
}

#8

}

\end{flushleft}

\end{minipage}


\vfill

\begin{center}
Vilnius\\
\the\year
\end{center}

\iftoggle{needPreface}{
 \sectionWithoutNumber{\keyWordPreface}{preface}
Pratarmės (Preface) informacija


\iftoggle{inLithuanian}
{
\vspace{\baselineskip}\hfill
\today
}
{
 \vspace{\baselineskip}\hfill \today
}

 \vspace{5cm}

\iftoggle{signaturesOnTitlePage}{}
{
\authorsignatures{#3}{#4}{#5}{#6}{#7}
}
}{}
\newpage
}


\begin{document}
% #1 -report type, #2 - title, #3-7 students, #8 - supervisor
\depttitlepage{Bachelors Thesis}{Implementation of application for visualization of regularities and randomness in data}{Audrius Baranauskas}
{}{}{}{}% students 2-5
{dr. Tadas Meškauskas}

\tableofcontents


%keywords and notations if needed
\sectionWithoutNumber{Keywords}{keywords}{Pateikiamas terminų sąrašas (jei reikia)}

%both abstracts
\bothabstracts{Santraukos tekstas rašto darbo kalba...}%tex-file of abstract in original language
{Darbo pavadinimas kita kalba} %if work is in LT this title should be in English
{Šio projekto tikslas buvo sukurti interneto aplikaciją, kuri sugebėti sukurti bei kategorizuoti rekurencines diagramas naudojamas duomenų duomenų dėsningumų ir atsitiktinų vizualizavimui.
Aplikacija sugeba reaguoti į varotojo įrenginio dydį.
Dėl šios priežasties ją galima naudotis tiek mobiliuoju telefonu, tiek namų kompiuteriu.
Duomenu kategorizacija yra įgyvendinta panaudojant kovoliucinius neuroninius tinklus.
Neuroninis tinklas gali tiksliai klasifikuoti chaotinius, periodinius bei trendą turinčius duomenis.
Ši publikacija detaliai aprašo aplikacijos vystymo procesą.
Ji taip pat nupasakoja kaip generuojami duomenys kovoliucinio neuroninio tinklo mokymams bei testavimui.
Galiausiai, buvo prieita išvada, kad konvoliucinis neuroninis tinklas gali išmokti kategorizuoti rekurentines diagramas. }%tex-file of abstract in other language


%Introduction section: label is sec:intro
\sectionWithoutNumber{\keyWordIntroduction}{intro}
Įvado tekstas ...


%=============================================================
%========================The main part========================
%=============================================================

\newpage

\section{Signal processing and Recurrence plot}
\label{sec:motivation}
\subsection{Signal processing}

A signal is a function that conveys information about the behaviour of a system
or attributes of some phenomenon \cite{priemer1990introductory}.
For example, measuring the time taken between a weight-driven
pendulum clock's ticks produces signal.
In turn, for the scope of this paper, we defined the term signal processing as
\textit{the science of analyzing time-varying processes} 
\cite{lyons2004understanding}.
By processing a signal we analyzed the non-triviality of a given signal.
Analyzing a signal reveals that some signals 
have properties that can be categorized.

\subsection{Signal property categories}


We have considered the following categories:
\begin{enumerate}
  \item Stationary and non stationary signals
  \item 
\end{enumerate}


Signals have varying properties. Some consist of simple repetitions while others
have no apparent patterns.
For example, measuring the time taken between a weight-driven
pendulum clock's ticks produces a relatively simple (trivial) signal.


% For the scope of this paper, we defined the term signal 
% processing as \textit{the science of analyzing time-varying 
% processes} \cite{lyons2004understanding}.




% \subsection{Signal classification}
% Signals have several classes. For the scope of this paper we considered
% the classes \cite{lathi1998signal}:
% \begin{enumerate}
%   \item Stationary and non-stationary signals
%   \item 
%   \item Periodic, quasi periodic and aperiodic signals 


% \end{enumerate}






% Some signals are inherently simple such as the time taken
% a falling object or

% \begin{lstlisting}
%   sygnalStates = []

%   # Generate data pairs, tripplets, quadruplets... D - plets
%   for i in range(0, self.M):
%       state = []

%       for j in range(0, self.D):
%           state.append(self.data[i+(j*self.d)])

%       sygnalStates.append(state)
% \end{lstlisting}


\newpage

\section{Pirmasis skyrius}
\label{sec:motivation}
\subsection{Pirmojo skyriaus poskyris}
\label{sec:example}
Pateikiamas \ref{sec:example} poskyrio tekstas. Vienas iš šaltinių~\cite{KTZ}. Visas ~\cite{KTV} turinys priklauso \ref{sec:motivation} skyriui.

\subsubsection{Pirmojo skyriaus pirmo poskyrio poskyris}
\label{sec:data}
Pateikiamas trečio lygio poskyrio tekstas.

\begin{equation}
  x = \sum_{i=1}^N m_i
\end{equation}

\begin{table}[!ht]\centering
  \caption{Lentelė ... }
  \label{tabl:table}
  \begin{tabular}{l|r|}
    test & test \\ \hline
    test & test \\
  \end{tabular}
\end{table}

Sprendimas pristatomas \ref{alg:1} algoritme, o įgyvendinimas -- \ref{abc} išeities kode.

\begin{algorithm}\caption{Algoritmas uždavinio sprendimui}
  \label{alg:1}
  \begin{algorithmic}
    \REQUIRE
    \ENSURE
    \STATE a \AND b
  \end{algorithmic}


\end{algorithm}



\begin{lstlisting}[caption={Pagrindinio metodo žingsniai},label={abc}]
public static void main(String args[]){
}
\end{lstlisting}

%Conclusions section
\sectionWithoutNumber{\keyWordConclusions}{conclu}
During the research and development phases of this project, the following goals were achieved:
\begin{itemize}
    \item The recurrence plot was analysed to evaluate what can be determinen from the visualization. Several observable characteristics were distinguished. A local implementation of the recurrence plot was developed and utilized throughout the application.
    \item A convolutional neural network, capable of identifying the characteristics of a recurrence plot, was developed and trained. The trained data model performed with high accuracy.
    \item An web application for generating recurrence plots was developed. A microservice architecture was used develop the application providing a modern, scalable solution. The application utilizes the aforementioned CNN to provide information about the generated, that a novel user might not distinguish.
\end{itemize}

The application functionality can be yet improved. Microservice could be containerized to provide easier deployment. The front end part of the application could support full CRUD operations for managing of data in the database. So far this this can only be done by accessing the backend service via an exposed API.

%ateities darbų gairės, planas/next steps of the work
\sectionWithoutNumber{Ateities tyrimų planas}{future}{Pristatomi ateities darbai ir/ar jų planas, gairės tolimesniems darbams....}


%file literatureSources.bib
\referenceSources{literatureSources}



%% this part is optional
\newpage
\begin{appendices}
  Dokumentą sudaro du priedai: \ref{app:a} priede  ....
  \newpage
  \section{Pirmojo priedo pavadinimas}
  \label{app:a}
  Pirmojo priedo tekstas ...

  \newpage
  \section{Antrojo priedo pavadinimas}
  Antrojo priedo tekstas ...

\end{appendices}


\end{document}
